\section{Scalar conductivities on a fixed Euclidean domain}
\label{scalar:section}

The geometric theorem determines a metric up to a change of coordinates.
For scalar conductivities on a fixed Euclidean domain, the corresponding
change of coordinates is conformal.  We first show that fixing a boundary
patch removes this remaining freedom, then verify the exact conversion
between the two measurement maps.

\begin{lemma}[Conformal rigidity at a boundary patch]
\label{scalar:boundary-rigidity}
Let $\Omega\subset\mathbb R^n$ be a connected open set with smooth
boundary, where $n\ge3$.  Suppose $F:\overline\Omega\to\overline\Omega$
is a smooth diffeomorphism through the boundary and
\[
 F^*e=e^{2b}e\quad\text{in }\Omega
\]
for a real $b\in C^\infty(\overline\Omega)$, where $e$ is the
Euclidean metric.  If $F$ fixes a nonempty relatively open subset of
$\partial\Omega$ pointwise, then $F=\Id$ and $b=0$.
\end{lemma}

\begin{proof}
The differential identities below are the classical local mechanism
behind Liouville's conformal rigidity theorem; see
\citet[Equations~(4)--(6)]{KushelmanMcGrath2024}.  We include the
argument with its boundary initial values, since preservation of the
interior side is essential here.

The Christoffel symbols of $e^{2b}e$ are
\[
 \Gamma^k_{ij}=b_i\delta^k_j+b_j\delta^k_i-b_k\delta_{ij},
 \qquad b_i=\partial_i b.
\]
Since $F$ is an isometry from this metric to the Euclidean metric,
the connection transformation rule gives
\begin{equation}\label{scalar:connection}
 D^2F(v,w)
 =DF\bigl(db(v)w+db(w)v-\langle v,w\rangle\nabla b\bigr).
\end{equation}
This identity initially holds in the interior and extends to the
boundary by smoothness.

Fix a point $q$ of the fixed boundary patch and let $\nu$ be the
Euclidean inward unit normal there.  For every tangent vector $v$,
differentiating $F=\Id$ along the fixed patch gives
$DF(q)v=v$.  Conformality therefore implies $b(q)=0$ and makes
$DF(q)$ orthogonal.  It follows that $DF(q)\nu$ is either $\nu$ or
$-\nu$.  To determine the sign, choose a smooth local defining
function $s$ positive in $\Omega$, with $ds(q)\nu=1$.  For small
$a>0$, both $q+a\nu$ and $F(q+a\nu)$ are interior points.  Hence
the right derivative at zero of $s(F(q+a\nu))$ is nonnegative.
It is $ds(q)DF(q)\nu=\pm1$, so the positive sign holds.
Consequently $DF(q)=I$ on the entire fixed patch.

The tangential derivatives of $b$ vanish there because $b=0$.
Choose a smooth curve $c$ in the patch with $c(0)=q$ and a nonzero
tangent $v=c'(0)$.  Twice differentiating $F(c(a))=c(a)$ gives
\[
 D^2F(q)(v,v)+DF(q)c''(0)=c''(0),
\]
so $D^2F(q)(v,v)=0$.  Equation~\eqref{scalar:connection}, together
with $DF(q)=I$ and $db(q)v=0$, now yields
$|v|^2\nabla b(q)=0$.  Thus both $b$ and $\nabla b$ vanish on the
fixed patch.  In particular, the argument rules out a conformal map
that fixes a hypersurface but exchanges its two sides.

The metric $e^{2b}e=F^*e$ is flat.  Substitution of the displayed
Christoffel symbols into its Ricci tensor gives
\[
 0=\operatorname{Ric}_{ij}
 =-(n-2)(b_{ij}-b_i b_j)
   -\bigl(\Delta b+(n-2)|\nabla b|^2\bigr)\delta_{ij}.
\]
Taking the Euclidean trace and substituting back, using $n-2\ne0$,
gives
\begin{equation}\label{scalar:hessian}
 b_{ij}=b_i b_j-\tfrac12|\nabla b|^2\delta_{ij}.
\end{equation}
Along a straight segment of constant velocity $v$, the vector
$q_b=\nabla b$ therefore solves
\[
 q_b'=(q_b\cdot v)q_b-\tfrac12|q_b|^2v.
\]
Start with a short inward segment at a point of the fixed patch.
Its initial value is zero, and the right side of this ordinary
differential equation is locally Lipschitz in $q_b$.  Uniqueness
implies $q_b=0$ along that segment.  Every connected open subset of
Euclidean space is polygonally connected: the set reachable from one
point by polygonal paths is both open and relatively closed, because
small balls lie in the domain.  Apply the same uniqueness argument
successively along those segments to conclude $\nabla b=0$ throughout
$\Omega$.  Continuity at the initial boundary point gives $b=0$.

Equation~\eqref{scalar:connection} now gives $D^2F=0$ in the connected
domain, so $F$ is affine there.  Its derivative and value at the fixed
boundary point are those of the identity.  Thus $F=\Id$ in $\Omega$
and on its closure by continuity.
\end{proof}

\begin{proof}[Proof of Corollary~\ref{main:scalar}]
For each conductivity define
\[
 g_j=\gamma_j^{\,2/(n-2)}e.
\]
These are smooth positive metrics on the compact manifold
$\overline\Omega$.  In Euclidean coordinates,
\begin{equation}\label{scalar:energy-conversion}
 \sqrt{\det g_j}\,g_j^{ab}=\gamma_j\delta^{ab},\qquad
 \int_\Omega\langle du,dv\rangle_{g_j}\,dV_{g_j}
       =\int_\Omega\gamma_j\nabla u\cdot\nabla v\,dx.
\end{equation}
Thus the $g_j$-harmonic functions are exactly the solutions of the
conductivity equation.  Their boundary energy output also has the
correct density: writing $\psi_j=\log\gamma_j/(n-2)$ gives
\[
 \nu_{g_j}=e^{-\psi_j}\nu_e,\qquad
 dS_{g_j}=e^{(n-1)\psi_j}dS_e,\qquad
 \partial_{\nu_{g_j}}u\,dS_{g_j}
       =\gamma_j\partial_{\nu_e}u\,dS_e.
\]
The equality of scalar local maps consequently gives equality of
the density-valued energy pairings for every smooth pair supported
in $\Gamma$, and hence of the maps on $H^{1/2}_{co}(\Gamma)$.

Choose a nonempty proper relatively open patch
$\Gamma_0\subset\Gamma$; if $\Gamma$ is already proper one may
take $\Gamma_0=\Gamma$.  Theorem~\ref{main:patch} gives a smooth
diffeomorphism $\Phi:\overline\Omega\to\overline\Omega$ with
$\Phi|_{\Gamma_0}=\Id$ and $g_2=\Phi^*g_1$.  Therefore
\[
 \Phi^*e=e^{2b}e,\qquad
 b=\frac{\log\gamma_2-\log\gamma_1\circ\Phi}{n-2}.
\]
Lemma~\ref{scalar:boundary-rigidity} gives $\Phi=\Id$.
The metric equality and positivity then imply
$\gamma_1=\gamma_2$ on $\overline\Omega$, as asserted.
\end{proof}

In dimension three, by contrast, two distinct uniformly positive bounded
measurable scalar conductivities on a ball, both equal to one near its
boundary, can have the same full weak Dirichlet-to-Neumann
operator~\cite[Theorem~1.1]{OpenAIRoughScalar2026}.
